In the next few chapters, we present a program analysis that can help data analysts design adaptive data analyses controlling their generalization errors. 
Given an input program implementing an adaptive data analysis, our program analysis generates an upper bound on the total number of queries that the data analysis will run, and more interestingly also an upper bound on the depth of the chain of queries. 
These two measures can be used to select the right technique to guarantee a bound on the generalization error of the data analysis. 
Our program analysis is based on an analysis of the dependency graph between different queries, representing the potential chain an adaptive data analysis may generate. 
We show how the proposed program analysis can help to analyze the generalization error of several concrete data analyses with different adaptivity.
\begin{figure}
    \centering
    \includegraphics[width=0.9\columnwidth]{chapters/adapt/overview-dissertation.png}
    \caption{Overview of our adaptive data analysis model. }
    \label{fig:adaptivity-model-overview}
\vspace{-0.5cm}
\end{figure}

